Status as
of
2026-10-05.
Scope: the
South
Korea
region of
the two
GCAM-KAIST
scenario
databases,
\texttt{KAIST\_\allowbreak{}9\_\allowbreak{}ref\_\allowbreak{}u0902v3}
(\texttt{ref})
and
\texttt{KAIST\_\allowbreak{}9\_\allowbreak{}NZ\_\allowbreak{}u0902v3}
(\texttt{nz}).
This is a
screening
audit.
Nothing in
it is a
reviewed
or
production-ready
input.

The
document
has two
parts:

\begin{itemize}
\tightlist
\item
  \textbf{Part
  A}
  explains,
  in plain
  language,
  what a
  GCAM-KAIST
  non-CO2
  number
  is made
  of and
  what we
  found.
\item
  \textbf{Part
  B} gives
  the
  technical
  detail
  and
  evidence
  behind
  every
  statement
  in Part
  A.
\end{itemize}

Where a
statement
is an
inference
rather
than
something
read
directly
from the
data, it
is marked
\textbf{Inferred}.

\hypertarget{part-a-plain-language-guide}{%
\subsection{Part
A:
Plain-language
guide}\label{part-a-plain-language-guide}}

\hypertarget{a1.-the-units-a-non-co2-number-is-built-from}{%
\subsubsection{A1.
The units
a non-CO2
number is
built
from}\label{a1.-the-units-a-non-co2-number-is-built-from}}

Every
non-CO2
emission
value
GCAM-KAIST
reports
belongs to
exactly
one
combination
of the
following:

\begin{longtable}[]{@{}
  >{\raggedright\arraybackslash}p{(\columnwidth - 4\tabcolsep) * \real{0.1481}}
  >{\raggedright\arraybackslash}p{(\columnwidth - 4\tabcolsep) * \real{0.4259}}
  >{\raggedright\arraybackslash}p{(\columnwidth - 4\tabcolsep) * \real{0.4259}}@{}}
\toprule\noalign{}
\begin{minipage}[b]{\linewidth}\raggedright
Unit
\end{minipage}
&
\begin{minipage}[b]{\linewidth}\raggedright
What it is
\end{minipage}
&
\begin{minipage}[b]{\linewidth}\raggedright
Example
\end{minipage} \\
\midrule\noalign{}
\endhead
\bottomrule\noalign{}
\endlastfoot
Scenario &
One model
run &
\texttt{ref},
\texttt{nz} \\
Region &
Geographic
unit &
South
Korea \\
Sector &
Something
GCAM
produces
or
supplies &
\texttt{electricity},
\texttt{process\ heat\ cement},
\texttt{trn\_\allowbreak{}pass\_\allowbreak{}road\_\allowbreak{}LDV\_\allowbreak{}4W}
(cars and
large
cars) \\
Subsector
& A group
inside a
sector,
usually by
fuel or
vehicle
type &
\texttt{gas},
\texttt{coal},
\texttt{Car} \\
Technology
& One
specific
way of
producing
the
sector's
output &
\texttt{gas\ (CC)}
(combined-cycle
gas
plant),
\texttt{Liquids}
(petrol/diesel
car) \\
Vintage &
The year a
piece of
equipment
was built
& the
2021-built
\texttt{gas\ (CC)}
fleet \\
Model year
& The year
being
reported &
2015,
2021,
2050 \\
Pollutant
& The
emitted
species &
\texttt{NOx},
\texttt{SO2\_\allowbreak{}2}
(sulfur
dioxide),
\texttt{NH3} \\
\end{longtable}

On top of
these
units sit
four
ideas.

\textbf{Activity.}
Every
technology
has
activity
in every
model
year: what
it
produces
(its
\emph{output},
for
example
exajoules
(EJ) of
electricity,
megatonnes
(Mt) of
cement, or
billions
of
passenger-km)
and what
it
consumes
(its
\emph{inputs},
for
example EJ
of coal).

\textbf{Non-CO2
tag.} A
non-CO2
emission
exists
only where
the model
attaches
an
emissions
object to
a
technology,
one per
pollutant.
In the
database
this is an
element
called
\texttt{Non-\allowbreak{}CO2}
(for
example
\texttt{\textless{}Non-\allowbreak{}CO2\ name=\allowbreak{}"NOx"\textgreater{}}).
In the
model
output the
tag holds
\textbf{only
one thing:
an
emission
mass for
each model
year}. It
does not
hold an
emission
factor,
the
activity
it is tied
to, or
anything
about
pollution-control
equipment.

\begin{quote}
If a
technology
has no tag
for a
pollutant,
GCAM-KAIST
reports
nothing
for it.
Not zero:
the
technology
simply
never
appears in
any
non-CO2
output.
\end{quote}

\textbf{Driver.}
Each tag
is tied to
one
activity
quantity,
called its
driver. In
the GCAM
input
files the
driver is
either a
named fuel
input
(coal
burned in
a cement
kiln,
petrol
burned by
cars) or
the
technology's
output
(electricity
generated,
tonnes of
steel,
tonnes of
meat).

\textbf{Emission
factor
(coefficient).}
The tag's
mass
divided by
its
driver:
for
example
kilograms
of NOx per
gigajoule
(GJ) of
coal
burned, or
per
megawatt-hour
(MWh)
generated.
\textbf{The
model
output
does not
store this
number.}
We
calculate
it, and
call it
the
\emph{implied}
coefficient.

Two more
details
matter
when
reading
the data.

\textbf{Vintages.}
In the
historical
years each
technology
has one
vintage.
In future
years
older
equipment
keeps
running
alongside
new
equipment,
and the
database
lists a
separate
tag value
for each
vintage.
For
example,
the NOx
tag on
coal
burned for
cement-kiln
heat has 3
vintage
entries in
2030 and 7
in 2050. A
technology's
total for
a year is
the sum
over its
vintages.

\textbf{Model
years.}
GCAM's
base
(calibration)
years are
1975,
1990,
2005,
2010, 2015
and 2021.
Future
model
years in
these runs
are 2025
to 2050 in
5-year
steps.
Technologies
are listed
out to
2100, but
carry no
output
after
2050.

\hypertarget{worked-example}{%
\paragraph{Worked
example}\label{worked-example}}

\texttt{ref},
2021,
technology
\texttt{coal}
in
subsector
\texttt{coal}
of sector
\texttt{process\ heat\ cement}
(coal
burned to
heat
cement
kilns):

\begin{longtable}[]{@{}
  >{\raggedright\arraybackslash}p{(\columnwidth - 2\tabcolsep) * \real{0.5541}}
  >{\raggedright\arraybackslash}p{(\columnwidth - 2\tabcolsep) * \real{0.4459}}@{}}
\toprule\noalign{}
\begin{minipage}[b]{\linewidth}\raggedright
Quantity
\end{minipage}
&
\begin{minipage}[b]{\linewidth}\raggedright
Value
\end{minipage} \\
\midrule\noalign{}
\endhead
\bottomrule\noalign{}
\endlastfoot
Driver:
coal input
(\texttt{delivered\ coal})
& 0.0835
EJ \\
NOx tag &
0.0220 Tg
(22.0
kt) \\
Implied
coefficient
= tag ÷
driver &
0.264 kg
NOx per GJ
of coal \\
\end{longtable}

\hypertarget{unit-conversions-used-in-this-document}{%
\paragraph{Unit
conversions
used in
this
document}\label{unit-conversions-used-in-this-document}}

\begin{longtable}[]{@{}
  >{\raggedright\arraybackslash}p{(\columnwidth - 2\tabcolsep) * \real{0.4762}}
  >{\raggedright\arraybackslash}p{(\columnwidth - 2\tabcolsep) * \real{0.5238}}@{}}
\toprule\noalign{}
\begin{minipage}[b]{\linewidth}\raggedright
From
\end{minipage}
&
\begin{minipage}[b]{\linewidth}\raggedright
To
\end{minipage} \\
\midrule\noalign{}
\endhead
\bottomrule\noalign{}
\endlastfoot
1 Tg
(teragram)
& 1,000 kt
= 1
million
tonnes \\
1 EJ
(exajoule)
& 10⁹ GJ =
2.778 ×
10⁸ MWh \\
Tg per EJ
& kg per
GJ \\
Tg per EJ
of
electricity
& × 3.6 →
kg per
MWh \\
Tg per Mt
& 1,000 kg
per
tonne \\
\end{longtable}

\hypertarget{a2.-how-a-number-gets-into-a-tag}{%
\subsubsection{A2.
How a
number
gets into
a
tag}\label{a2.-how-a-number-gets-into-a-tag}}

There are
two
phases.

\textbf{Base
years (up
to 2021):
the mass
is given.}
For most
sectors,
GCAM's
input data
supplies a
total
emission
mass for
each
technology,
pollutant
and year.
GCAM then
divides
that mass
by its own
modelled
activity
for the
technology
to get the
coefficient.
The
coefficient
is not
measured
or looked
up; it is
whatever
is left
after the
division.

For a few
sectors
(iron and
steel,
fuel
extraction,
and two
crop
residue-burning
species)
the input
data gives
the
coefficient
directly
instead.

\textbf{Future
years
(2025 to
2050):
emissions
=
coefficient
×
activity.}
The
coefficient
is either
kept at
its 2021
value or
reduced
along a
pre-set
curve. In
the stock
GCAM
inputs,
that curve
is tied to
GDP.
Almost all
coefficients
are the
same in
both
scenarios,
so what
differs
between
\texttt{ref}
and
\texttt{nz}
is mainly
the
activity.

\hypertarget{a3.-what-we-found}{%
\subsubsection{A3.
What we
found}\label{a3.-what-we-found}}

\begin{enumerate}
\def\labelenumi{\arabic{enumi}.}
\tightlist
\item
  \textbf{Only
  tagged
  technologies
  report,
  and some
  large
  activities
  have no
  tags.}
  All
  fossil
  power
  plants
  in
  operation
  are
  tagged.
  But
  several
  large
  activities
  are not:

  \begin{itemize}
  \tightlist
  \item
    \textbf{Cement
    making.}
    Only
    the
    fuel
    burned
    for
    kiln
    heat
    is
    tagged;
    the
    kiln
    process
    itself
    is
    not.
    Paper,
    food
    processing
    and
    chemicals
    are
    the
    same:
    fuel
    only.
  \item
    \textbf{Oil
    refining.}
    No tag
    at
    all,
    despite
    4.69
    EJ of
    output
    in
    2021.
  \item
    \textbf{Chemical
    feedstocks
    and
    industrial
    cogeneration.}
  \item
    \textbf{Several
    technologies
    that
    grow
    in the
    net-zero
    scenario:}
    hybrid
    vehicles,
    hydrogen-based
    steel,
    and
    steel
    and
    cement
    with
    carbon
    capture.
    These
    show
    zero
    emissions
    even
    though,
    for
    example,
    hydrogen-based
    steel
    still
    burns
    gas
    and
    oil.
  \end{itemize}
\item
  \textbf{For
  most
  sectors,
  base-year
  numbers
  are
  inventory
  numbers,
  not
  model
  results.}

  \begin{itemize}
  \tightlist
  \item
    GCAM-KAIST's
    2015
    and
    2021
    tags
    equal
    those
    in
    stock
    GCAM
    v9
    (2,244
    of
    2,251
    checked
    values
    match
    within
    0.1\%),
    so
    KAIST
    largely
    kept
    GCAM's
    standard
    calibration.
  \item
    Stock
    GCAM's
    calibration
    comes
    from
    CEDS,
    the
    Community
    Emissions
    Data
    System.
    CEDS
    documents
    scaling
    its
    Korean
    estimates
    to
    NIER's
    national
    inventory,
    which
    is
    published
    as
    CAPSS.
  \item
    Eleven
    GCAM-KAIST
    totals
    reproduce
    CAPSS
    category
    totals
    to
    within
    0.12\%,
    for
    example
    agricultural
    NH3:
    200,384.0
    t vs
    CAPSS
    200,384
    t in
    2021.
  \end{itemize}
\item
  \textbf{Technologies
  inside
  the same
  inventory
  total
  share
  one
  coefficient.}
  Because
  the
  coefficient
  is a
  by-product
  of
  dividing
  a total,
  distinct
  technologies
  end up
  with the
  same
  value:

  \begin{itemize}
  \tightlist
  \item
    combined-cycle
    and
    steam/turbine
    gas
    plants
    have
    identical
    coefficients
    per
    unit
    of
    fuel
    in
    every
    base
    year
    from
    2005
    to
    2021;
  \item
    commercial
    heating,
    cooling
    and
    ``other''
    uses
    of gas
    share
    one
    coefficient;
  \item
    all
    ten
    household
    income
    groups
    share
    one
    coefficient;
  \item
    mobile
    and
    stationary
    machinery
    in
    agriculture,
    construction
    and
    mining
    carry
    exactly
    equal
    masses
    in
    every
    base
    year
    from
    2005
    to
    2021.
  \end{itemize}
\item
  \textbf{Some
  emissions
  are not
  tied to
  any
  activity.}
  Industrial
  process,
  solvent
  and
  waste
  emissions
  are
  attached
  to a
  fixed
  placeholder,
  not to
  production.
  In 2021
  these
  placeholder-driven
  tags
  hold
  81\% of
  GCAM-KAIST's
  Korean
  VOCs and
  55\% of
  its SO2,
  excluding
  international
  shipping
  and
  aviation.
\item
  \textbf{Base-year
  coefficients
  are
  unstable.}
  For the
  typical
  technology
  and
  pollutant,
  the
  implied
  coefficient
  varies
  by a
  factor
  of 2
  between
  2005 and
  2021;
  for one
  in ten
  it
  varies
  by more
  than
  11×.
\item
  \textbf{Future
  coefficients
  don't
  respond
  to
  policy.}
  From
  2021 to
  2050:

  \begin{itemize}
  \tightlist
  \item
    46\%
    of
    coefficients
    stay
    frozen
    at
    their
    2021
    value;
  \item
    52\%
    decline
    along
    the
    pre-set
    curve;
  \item
    in
    2050,
    93\%
    are
    identical
    in
    \texttt{ref}
    and
    \texttt{nz}.
  \end{itemize}
\item
  \textbf{No
  primary
  particulate
  matter.}
  GCAM-KAIST
  reports
  black
  carbon
  (BC) and
  organic
  carbon
  (OC),
  but not
  PM2.5,
  PM10 or
  total
  suspended
  particulates.
\end{enumerate}

\hypertarget{a4.-glossary}{%
\subsubsection{A4.
Glossary}\label{a4.-glossary}}

\begin{longtable}[]{@{}
  >{\raggedright\arraybackslash}p{(\columnwidth - 2\tabcolsep) * \real{0.4250}}
  >{\raggedright\arraybackslash}p{(\columnwidth - 2\tabcolsep) * \real{0.5750}}@{}}
\toprule\noalign{}
\begin{minipage}[b]{\linewidth}\raggedright
Term
\end{minipage}
&
\begin{minipage}[b]{\linewidth}\raggedright
Meaning
\end{minipage} \\
\midrule\noalign{}
\endhead
\bottomrule\noalign{}
\endlastfoot
GCAM &
Global
Change
Analysis
Model
(PNNL/JGCRI).
GCAM-KAIST
is KAIST's
version of
it, run
here with
the
scenario
names
above. \\
Stock GCAM
v9 & The
standard
JGCRI
release,
unpacked
locally at
\texttt{.\allowbreak{}.\allowbreak{}/\allowbreak{}gcam-\allowbreak{}v9},
used as
the
comparison
baseline. \\
CEDS &
Community
Emissions
Data
System. A
global
historical
emissions
inventory;
GCAM
calibrates
its
base-year
non-CO2
emissions
to it. \\
CAPSS &
Clean Air
Policy
Support
System.
Korea's
national
air
pollutant
emissions
inventory,
published
by NIER
(National
Institute
of
Environmental
Research). \\
Base year
/
calibration
year & A
historical
model year
in which
GCAM is
fitted to
observed
data
(1975--2021
here). \\
Pass-through
sector & A
GCAM
sector
that only
routes a
flow, such
as the
\texttt{elec\_\allowbreak{}*}
sectors
that
choose
power-plant
cooling
systems. \\
\texttt{\_\allowbreak{}AGR},
\texttt{\_\allowbreak{}AWB}
&
Pollutant-name
suffixes:
agricultural
emissions
(fertilizer,
livestock)
and
agricultural
waste
(crop
residue)
burning. \\
\texttt{SO2\_\allowbreak{}2}
& GCAM's
name for
sulfur
dioxide. \\
NMVOC /
VOCs &
Non-methane
volatile
organic
compounds.
CAPSS
calls them
VOCs. \\
BaseX &
The XML
database
engine the
scenario
outputs
are stored
in. \\
\end{longtable}

\hypertarget{part-b-technical-audit}{%
\subsection{Part
B:
Technical
audit}\label{part-b-technical-audit}}

\hypertarget{b1.-data-tools-and-definitions}{%
\subsubsection{B1.
Data,
tools and
definitions}\label{b1.-data-tools-and-definitions}}

\textbf{GCAM-KAIST
outputs.}

\begin{itemize}
\tightlist
\item
  Two
  BaseX
  databases
  under
  \texttt{model\_\allowbreak{}inputs/\allowbreak{}gcam/\allowbreak{}gcam\_\allowbreak{}kaist/\allowbreak{}},
  registered
  as
  \texttt{ref}
  and
  \texttt{nz}.
\item
  Each
  holds
  one
  scenario:
  \texttt{KAIST\_\allowbreak{}9\_\allowbreak{}ref\_\allowbreak{}u0902v3}
  and
  \texttt{KAIST\_\allowbreak{}9\_\allowbreak{}NZ\_\allowbreak{}u0902v3},
  with
  \texttt{date}
  attributes
  2026-02-09T14:37:41+09:00
  and
  14:39:03+09:00.
\item
  Only the
  region
  \texttt{South\ Korea}
  is
  analysed.
\end{itemize}

\textbf{GCAM-KAIST
inputs.}
The input
XMLs KAIST
used to
build
these runs
are
\textbf{not
available
locally}.
Every
statement
about how
GCAM-KAIST
sets its
tags is
therefore
based on
its
outputs,
compared
with stock
GCAM v9's
inputs.

\textbf{Stock
GCAM v9.}
\texttt{.\allowbreak{}.\allowbreak{}/\allowbreak{}gcam-\allowbreak{}v9},
the JGCRI
GCAM
release.
Its
\texttt{input/\allowbreak{}gcamdata}
R code and
\texttt{extdata}
contain no
reference
to CAPSS
or KAIST.
Files
used:

\begin{longtable}[]{@{}
  >{\raggedright\arraybackslash}p{(\columnwidth - 2\tabcolsep) * \real{0.5412}}
  >{\raggedright\arraybackslash}p{(\columnwidth - 2\tabcolsep) * \real{0.4588}}@{}}
\toprule\noalign{}
\begin{minipage}[b]{\linewidth}\raggedright
File
\end{minipage}
&
\begin{minipage}[b]{\linewidth}\raggedright
Used for
\end{minipage} \\
\midrule\noalign{}
\endhead
\bottomrule\noalign{}
\endlastfoot
\texttt{input/\allowbreak{}gcamdata/\allowbreak{}xml/\allowbreak{}all\_\allowbreak{}energy\_\allowbreak{}emissions.\allowbreak{}xml}
&
Energy-sector
base-year
tags \\
\texttt{input/\allowbreak{}gcamdata/\allowbreak{}xml/\allowbreak{}all\_\allowbreak{}aglu\_\allowbreak{}emissions\_\allowbreak{}IRR\_\allowbreak{}MGMT.\allowbreak{}xml}
&
Agriculture
tags \\
\texttt{input/\allowbreak{}gcamdata/\allowbreak{}xml/\allowbreak{}all\_\allowbreak{}unmgd\_\allowbreak{}emissions.\allowbreak{}xml},
\texttt{all\_\allowbreak{}protected\_\allowbreak{}unmgd\_\allowbreak{}emissions.\allowbreak{}xml}
&
Unmanaged-land
tags \\
\texttt{input/\allowbreak{}gcamdata/\allowbreak{}xml/\allowbreak{}emission\_\allowbreak{}factor\_\allowbreak{}controls.\allowbreak{}xml}
&
Future-year
coefficient
overrides \\
\texttt{output/\allowbreak{}queries/\allowbreak{}Main\_\allowbreak{}queries.\allowbreak{}xml}
& Stock
reporting
queries \\
\texttt{input/\allowbreak{}gcamdata/\allowbreak{}R/\allowbreak{}constants.\allowbreak{}R}
&
Base-year
definition \\
\texttt{input/\allowbreak{}gcamdata/\allowbreak{}inst/\allowbreak{}extdata/\allowbreak{}emissions/\allowbreak{}CEDS/\allowbreak{}ceds\_\allowbreak{}sector\_\allowbreak{}map.\allowbreak{}csv}
& CEDS →
GCAM
sector
mapping \\
\end{longtable}

\textbf{CAPSS.}
\texttt{docs/\allowbreak{}2015\%20Emissions\%20by\%20Source\%20Category.\allowbreak{}xlsx}
and
\texttt{docs/\allowbreak{}2021\%20Emissions\%20by\%20Source\%20Category.\allowbreak{}xlsx}.
These hold
first-level
source-category
totals
only, in
tonnes.

\textbf{Definitions.}

\begin{itemize}
\tightlist
\item
  \textbf{Active:}
  a
  technology
  with
  output
  above
  zero in
  at least
  one
  model
  year up
  to 2050,
  in
  either
  scenario.
\item
  \textbf{Air-quality
  (AQ)
  species:}
  BC, OC,
  CO, NH3,
  NMVOC,
  NOx, SO2
  (including
  their
  \texttt{\_\allowbreak{}AGR}
  and
  \texttt{\_\allowbreak{}AWB}
  variants
  unless
  stated).
\item
  \textbf{Identical:}
  within
  0.1\%
  relative
  difference,
  unless
  stated
  otherwise.
\end{itemize}

\hypertarget{b2.-how-tags-are-stored}{%
\subsubsection{B2.
How tags
are
stored}\label{b2.-how-tags-are-stored}}

\textbf{Element
structure.}
A tag is a
child of a
technology
element:

\begin{Verbatim}[breaklines,fontsize=\footnotesize]
<technology year="2021" name="other industrial processes" type="technology">
  <output-primary name="industrial processes" type="output">
    <physical-output unit="NA" vintage="2021">0.000243714</physical-output>
  </output-primary>
  <!-- share-weight and the BC, CH4, CO, N2O, NH3, NMVOC tags omitted -->
  <Non-CO2 name="NOx" type="GHG">
    <emissions unit="Tg" year="2021">0.0390991</emissions>
  </Non-CO2>
  <input-energy name="misc emissions sources" type="input">
    <demand-physical unit="NA" vintage="2021">0.000243714</demand-physical>
    ...
\end{Verbatim}

\begin{itemize}
\tightlist
\item
  In the
  output
  database,
  the only
  child
  element
  of any
  \texttt{Non-\allowbreak{}CO2}
  is
  \texttt{emissions}.
  No
  coefficient,
  driver,
  control
  parameter
  or data
  source
  is
  stored.
\item
  CO2 is a
  separate
  element
  named
  \texttt{CO2}.
  Every
  element
  with
  \texttt{type=\allowbreak{}"GHG"}
  is
  either
  \texttt{CO2}
  or
  \texttt{Non-\allowbreak{}CO2}.
  In
  \texttt{ref},
  South
  Korea
  has
  23,117
  \texttt{Non-\allowbreak{}CO2}
  and
  1,676
  \texttt{CO2}
  elements,
  counting
  each
  technology
  vintage
  separately.
\end{itemize}

\textbf{Where
tags
occur.}
\texttt{Non-\allowbreak{}CO2}
appears
under
exactly
six
element
paths:

\begin{longtable}[]{@{}
  >{\raggedright\arraybackslash}p{(\columnwidth - 0\tabcolsep) * \real{1.0000}}@{}}
\toprule\noalign{}
\begin{minipage}[b]{\linewidth}\raggedright
Path
(below
\texttt{region})
\end{minipage} \\
\midrule\noalign{}
\endhead
\bottomrule\noalign{}
\endlastfoot
\texttt{supplysector/\allowbreak{}subsector/\allowbreak{}technology/\allowbreak{}Non-\allowbreak{}CO2} \\
\texttt{supplysector/\allowbreak{}subsector/\allowbreak{}pass-\allowbreak{}through-\allowbreak{}technology/\allowbreak{}Non-\allowbreak{}CO2} \\
\texttt{supplysector/\allowbreak{}tranSubsector/\allowbreak{}tranTechnology/\allowbreak{}Non-\allowbreak{}CO2} \\
\texttt{AgSupplySector/\allowbreak{}AgSupplySubsector/\allowbreak{}AgProductionTechnology/\allowbreak{}Non-\allowbreak{}CO2} \\
\texttt{AgSupplySector/\allowbreak{}AgSupplySubsector/\allowbreak{}UnmanagedLandTechnology/\allowbreak{}Non-\allowbreak{}CO2} \\
\texttt{resource/\allowbreak{}reserve-\allowbreak{}subresource/\allowbreak{}resource-\allowbreak{}reserve-\allowbreak{}technology/\allowbreak{}Non-\allowbreak{}CO2} \\
\end{longtable}

It never
appears
under
\texttt{pass-\allowbreak{}through-\allowbreak{}sector},
\texttt{renewresource},
\texttt{unlimited-\allowbreak{}resource},
or
technologies
nested
under a
\texttt{nesting-\allowbreak{}subsector}.

\textbf{Units
and
species}
(South
Korea,
\texttt{ref}):

\begin{longtable}[]{@{}
  >{\raggedright\arraybackslash}p{(\columnwidth - 2\tabcolsep) * \real{0.3134}}
  >{\raggedright\arraybackslash}p{(\columnwidth - 2\tabcolsep) * \real{0.6866}}@{}}
\toprule\noalign{}
\begin{minipage}[b]{\linewidth}\raggedright
Unit
\end{minipage}
&
\begin{minipage}[b]{\linewidth}\raggedright
Species
\end{minipage} \\
\midrule\noalign{}
\endhead
\bottomrule\noalign{}
\endlastfoot
Tg & BC,
CH4, CO,
H2, N2O,
NH3,
NMVOC,
NOx, OC,
SO2\_2;
CH4\_AGR,
N2O\_AGR,
NH3\_AGR,
NMVOC\_AGR,
NOx\_AGR;
BC\_AWB,
CH4\_AWB,
CO\_AWB,
H2\_AWB,
N2O\_AWB,
NH3\_AWB,
NMVOC\_AWB,
NOx\_AWB,
OC\_AWB,
SO2\_2\_AWB \\
Gg & C2F6,
CF4,
HFC125,
HFC134a,
HFC143a,
HFC227ea,
HFC23,
HFC236fa,
HFC32,
HFC43,
SF6 \\
MTC (Mt
carbon) &
CO2\_FUG
(fugitive
CO2,
stored
under
\texttt{Non-\allowbreak{}CO2}) \\
\end{longtable}

There is
no PM2.5,
PM10 or
TSP
species.

\textbf{Years
and
vintages.}

\begin{itemize}
\tightlist
\item
  Technology
  elements
  repeat
  once per
  vintage,
  using
  the
  \texttt{year}
  attribute.
\item
  Within a
  technology,
  \texttt{emissions/\allowbreak{}@year}
  and
  \texttt{physical-\allowbreak{}output/\allowbreak{}@vintage}
  both
  give the
  \emph{model
  year}
  being
  reported,
  despite
  the
  attribute
  name.
\item
  Base
  years
  have one
  vintage
  per
  technology.
  Future
  years
  have
  several
  (cement-kiln
  coal
  NOx: 3
  in 2030,
  7 in
  2050).
\item
  Model
  years
  present:
  1975,
  1990,
  2005,
  2010,
  2015,
  2021,
  2025,
  2030,
  2035,
  2040,
  2045,
  2050.
  Technology
  elements
  for
  2055--2100
  carry
  only
  share
  weights.
\item
  Stock
  \texttt{constants.\allowbreak{}R}
  defines
  \texttt{MODEL\_\allowbreak{}BASE\_\allowbreak{}YEARS}
  as 1975,
  1990,
  2005,
  2010,
  2015 and
  the
  final
  historical
  year,
  which is
  2021 in
  these
  outputs.
\end{itemize}

\textbf{Scenario
independence
of base
years.}
For model
years up
to 2021,
11,519 of
11,521
technology-pollutant-year
emission
values are
exactly
equal in
\texttt{nz}
and
\texttt{ref}.
The other
two differ
by at most
1 × 10⁻⁹
Tg.

\hypertarget{b3.-how-the-model-reports-them-native-queries}{%
\subsubsection{B3.
How the
model
reports
them:
native
queries}\label{b3.-how-the-model-reports-them-native-queries}}

Stock
GCAM's
\texttt{output/\allowbreak{}queries/\allowbreak{}Main\_\allowbreak{}queries.\allowbreak{}xml}
defines
five
non-CO2
queries:

\begin{longtable}[]{@{}
  >{\raggedright\arraybackslash}p{(\columnwidth - 2\tabcolsep) * \real{0.1786}}
  >{\raggedright\arraybackslash}p{(\columnwidth - 2\tabcolsep) * \real{0.8214}}@{}}
\toprule\noalign{}
\begin{minipage}[b]{\linewidth}\raggedright
Line
\end{minipage}
&
\begin{minipage}[b]{\linewidth}\raggedright
Query
title
\end{minipage} \\
\midrule\noalign{}
\endhead
\bottomrule\noalign{}
\endlastfoot
4916 &
\texttt{nonCO2\ emissions\ by\ region} \\
4928 &
\texttt{nonCO2\ emissions\ by\ sector\ (excluding\ resource\ production)} \\
4934 &
\texttt{nonCO2\ emissions\ by\ resource\ production} \\
4945 &
\texttt{nonCO2\ emissions\ by\ subsector\ (excluding\ resource\ production)} \\
4952 &
\texttt{nonCO2\ emissions\ by\ tech\ (excluding\ resource\ production)} \\
\end{longtable}

Each
selects
\texttt{*{[}@type\ =\allowbreak{}\ \textquotesingle{}GHG\textquotesingle{}\ and\ @name\ !=\allowbreak{}\ \textquotesingle{}CO2\textquotesingle{}\ and\ @name\ !=\allowbreak{}\ \textquotesingle{}CO2\_\allowbreak{}FUG\textquotesingle{}{]}/\allowbreak{}emissions}
beneath
sector
and/or
resource
nodes.
None of
them names
a sector.

\textbf{Consequence:
there are
no
sector-specific
non-CO2
queries.}
A sector
reports
non-CO2
emissions
if, and
only if,
its
technologies
carry
\texttt{Non-\allowbreak{}CO2}
tags. ``Is
there a
query for
cement
non-CO2?''
reduces to
``does any
cement
technology
carry a
tag?'',
and the
answer is
no (B4).

The repo's
query
\texttt{src/\allowbreak{}gcam\_\allowbreak{}emissions/\allowbreak{}gcam/\allowbreak{}queries/\allowbreak{}nonco2\_\allowbreak{}emissions\_\allowbreak{}by\_\allowbreak{}sector.\allowbreak{}xq}
selects
every
\texttt{Non-\allowbreak{}CO2/\allowbreak{}emissions}.
That is
the same
set of
tags as
the stock
queries,
plus
\texttt{CO2\_\allowbreak{}FUG}.

\hypertarget{b4.-coverage-which-activity-carries-tags}{%
\subsubsection{B4.
Coverage:
which
activity
carries
tags}\label{b4.-coverage-which-activity-carries-tags}}

\textbf{Method.}
\texttt{results/\allowbreak{}diagnostics/\allowbreak{}nonco2\_\allowbreak{}coverage/\allowbreak{}activity\_\allowbreak{}inventory.\allowbreak{}xq}
lists
every
technology
under
\texttt{supplysector},
\texttt{AgSupplySector},
\texttt{pass-\allowbreak{}through-\allowbreak{}sector},
\texttt{resource},
\texttt{renewresource}
and
\texttt{unlimited-\allowbreak{}resource}.
For each
it records
output by
year and
the
\texttt{Non-\allowbreak{}CO2}
species it
carries.
The set of
tagged
(sector,
subsector,
technology)
triples it
finds, 286
in total,
is
identical
to the set
in the
repo's
emissions
detail
table.

\textbf{Counts}
(active
nodes
only; each
household
income
decile
counts as
its own
sector):

\begin{longtable}[]{@{}
  >{\raggedright\arraybackslash}p{(\columnwidth - 4\tabcolsep) * \real{0.5606}}
  >{\raggedright\arraybackslash}p{(\columnwidth - 4\tabcolsep) * \real{0.1970}}
  >{\raggedright\arraybackslash}p{(\columnwidth - 4\tabcolsep) * \real{0.2424}}@{}}
\toprule\noalign{}
\begin{minipage}[b]{\linewidth}\raggedright
\end{minipage}
&
\begin{minipage}[b]{\linewidth}\raggedright
Sectors
\end{minipage}
&
\begin{minipage}[b]{\linewidth}\raggedright
Subsectors
\end{minipage} \\
\midrule\noalign{}
\endhead
\bottomrule\noalign{}
\endlastfoot
All active
technologies
tagged &
77 &
169 \\
Some
active
technologies
tagged &
50 & 36 \\
No active
technology
tagged &
139 &
281 \\
\textbf{Total
active} &
\textbf{266}
&
\textbf{486} \\
\end{longtable}

No tagged
technology
has
all-zero
emissions
in every
year up to
2050.

\textbf{Power.}

\begin{itemize}
\tightlist
\item
  Power
  tags sit
  on
  \texttt{electricity/\allowbreak{}\textless{}fuel\textgreater{}/\allowbreak{}\textless{}pass-\allowbreak{}through-\allowbreak{}technology\textgreater{}},
  for
  example
  \texttt{electricity/\allowbreak{}gas/\allowbreak{}gas\ (CC)}.
\item
  That
  technology's
  only
  energy
  input is
  the
  pass-through
  sector
  \texttt{elec\_\allowbreak{}gas\ (CC)},
  in a
  quantity
  equal to
  the
  electricity
  generated.
\item
  The 18
  \texttt{elec\_\allowbreak{}*}
  pass-through
  sectors,
  which
  hold the
  cooling-system
  choices,
  carry no
  \texttt{Non-\allowbreak{}CO2}
  tags.
  The
  fossil
  ones
  carry
  \texttt{CO2}
  tags.
\item
  Every
  fossil
  power
  technology
  with
  activity
  is
  tagged.
  The
  untagged
  fossil
  power
  technologies
  have
  zero
  output
  through
  2050 in
  both
  scenarios:
  \texttt{coal\ (IGCC)},
  \texttt{coal\ (conv\ pul\ CCS)},
  \texttt{coal\ (IGCC\ CCS)},
  \texttt{refined\ liquids\ (CC)},
  \texttt{refined\ liquids\ (CC\ CCS)},
  \texttt{biomass\ (conv\ CCS)}
  and
  \texttt{biomass\ (IGCC\ CCS)}.
\end{itemize}

\textbf{Active
technologies
with no
tag that
burn fuel
on site or
run an
industrial
process:}

\begin{longtable}[]{@{}
  >{\raggedright\arraybackslash}p{(\columnwidth - 10\tabcolsep) * \real{0.3622}}
  >{\raggedright\arraybackslash}p{(\columnwidth - 10\tabcolsep) * \real{0.1339}}
  >{\raggedright\arraybackslash}p{(\columnwidth - 10\tabcolsep) * \real{0.1102}}
  >{\raggedright\arraybackslash}p{(\columnwidth - 10\tabcolsep) * \real{0.1102}}
  >{\raggedright\arraybackslash}p{(\columnwidth - 10\tabcolsep) * \real{0.1024}}
  >{\raggedright\arraybackslash}p{(\columnwidth - 10\tabcolsep) * \real{0.1811}}@{}}
\toprule\noalign{}
\begin{minipage}[b]{\linewidth}\raggedright
Sector /
technology
\end{minipage}
&
\begin{minipage}[b]{\linewidth}\raggedright
Output
unit
\end{minipage}
&
\begin{minipage}[b]{\linewidth}\raggedright
2021
\texttt{ref}
\end{minipage}
&
\begin{minipage}[b]{\linewidth}\raggedright
2050
\texttt{ref}
\end{minipage}
&
\begin{minipage}[b]{\linewidth}\raggedright
2050
\texttt{nz}
\end{minipage}
&
\begin{minipage}[b]{\linewidth}\raggedright
Carries a
CO2 tag
\end{minipage} \\
\midrule\noalign{}
\endhead
\bottomrule\noalign{}
\endlastfoot
\texttt{cement}
/
\texttt{cement}
& Mt &
50.45 &
47.87 &
6.774 &
yes \\
\texttt{cement}
/
\texttt{cement\ CCS}
& Mt & 0 &
0 & 30.41
& yes \\
\texttt{paper}
/
\texttt{paper}
& Mt &
10.83 &
9.917 &
9.741 &
no \\
\texttt{food\ processing}
/
\texttt{food\ processing}
& Pcal &
59.36 &
56.43 &
56.22 &
no \\
\texttt{chemical}
/
\texttt{chemical}
& EJ &
2.422 &
3.150 &
3.026 &
no \\
\texttt{chemical\ feedstocks}
/
\texttt{refined\ liquids}
& EJ &
2.086 &
2.453 &
2.408 &
yes \\
\texttt{other\ industry}
/
\texttt{other\ industry}
& EJ &
1.090 &
1.314 &
1.289 &
no \\
\texttt{construction}
/
\texttt{construction}
& EJ &
0.0871 &
0.0850 &
0.0810 &
no \\
\texttt{N\ fertilizer}
/
\texttt{N\ fertilizer}
& Mt N &
0.2279 &
0.2058 &
0.1757 &
no \\
\texttt{refining}
/
\texttt{oil\ refining}
& EJ &
4.694 &
3.666 &
2.405 &
yes \\
\texttt{other\ industrial\ energy\ use}
/
\texttt{coal\ cogen}
& EJ &
0.1436 &
0.06497 &
0.00004 &
yes \\
\texttt{other\ industrial\ energy\ use}
/
\texttt{gas\ cogen}
& EJ &
0.0473 &
0.01101 &
0.0003 &
yes \\
\texttt{other\ industrial\ energy\ use}
/
\texttt{refined\ liquids\ cogen}
& EJ &
0.0408 &
0.00411 &
0.00002 &
yes \\
\texttt{other\ industrial\ energy\ use}
/
\texttt{biomass\ cogen}
& EJ &
0.00374 &
0.00228 &
0.00008 &
yes \\
\texttt{iron\ and\ steel}
/
\texttt{Hydrogen-\allowbreak{}based\ DRI}
& Mt & 0 &
10.29 &
16.36 &
yes \\
\texttt{iron\ and\ steel}
/
\texttt{BLASTFUR\ CCS}
& Mt & 0 &
0 & 4.150
& yes \\
\end{longtable}

None of the technologies in this table has an entry in stock GCAM v9's
Korea block of
\texttt{all\_\allowbreak{}energy\_\allowbreak{}emissions.\allowbreak{}xml}.
\texttt{Hybrid\ Liquids} and \texttt{oil\ refining} have no entry for any
region. These gaps are therefore inherited from stock GCAM, not introduced
by KAIST.

The
fuel-combustion
emissions
of
\texttt{cement},
\texttt{paper},
\texttt{food\ processing},
\texttt{chemical},
\texttt{other\ industry}
and
\texttt{construction}
\emph{are}
tagged,
but on
companion
energy-use
sectors:
\texttt{process\ heat\ cement},
\texttt{process\ heat\ paper},
\texttt{waste\ biomass\ for\ paper},
\texttt{process\ heat\ food\ processing},
\texttt{chemical\ energy\ use},
\texttt{other\ industrial\ energy\ use}
and
\texttt{construction\ energy\ use}.
What has
no tag
anywhere
is the
non-combustion
process
emissions
of these
industries,
apart from
the
unattributed
lump in
B6.4.

\textbf{Untagged
technologies
still burn
fuel.} In
\texttt{nz}
2050:

\begin{itemize}
\tightlist
\item
  \texttt{Hydrogen-\allowbreak{}based\ DRI}
  consumes
  0.0428
  EJ of
  \texttt{wholesale\ gas}
  and
  0.0180
  EJ of
  \texttt{refined\ liquids\ industrial},
  as well
  as
  hydrogen
  and
  electricity.
\item
  \texttt{BLASTFUR\ CCS}
  consumes
  0.0617
  EJ of
  \texttt{delivered\ coal}
  and
  0.0238
  EJ of
  \texttt{wholesale\ gas}.
\end{itemize}

\textbf{Hybrid
vehicles}
(\texttt{Hybrid\ Liquids})
have no
tags in
any mode:

\begin{longtable}[]{@{}
  >{\raggedright\arraybackslash}p{(\columnwidth - 6\tabcolsep) * \real{0.4646}}
  >{\raggedright\arraybackslash}p{(\columnwidth - 6\tabcolsep) * \real{0.2626}}
  >{\raggedright\arraybackslash}p{(\columnwidth - 6\tabcolsep) * \real{0.1414}}
  >{\raggedright\arraybackslash}p{(\columnwidth - 6\tabcolsep) * \real{0.1313}}@{}}
\toprule\noalign{}
\begin{minipage}[b]{\linewidth}\raggedright
Sector /
subsector
\end{minipage}
&
\begin{minipage}[b]{\linewidth}\raggedright
Unit
\end{minipage}
&
\begin{minipage}[b]{\linewidth}\raggedright
2050
\texttt{ref}
\end{minipage}
&
\begin{minipage}[b]{\linewidth}\raggedright
2050
\texttt{nz}
\end{minipage} \\
\midrule\noalign{}
\endhead
\bottomrule\noalign{}
\endlastfoot
\texttt{trn\_\allowbreak{}pass\_\allowbreak{}road}
/
\texttt{Bus}
& billion
passenger-km
& 108.9 &
78.02 \\
\texttt{trn\_\allowbreak{}pass\_\allowbreak{}road\_\allowbreak{}LDV\_\allowbreak{}4W}
/
\texttt{Car}
& billion
passenger-km
& 42.91 &
19.16 \\
\texttt{trn\_\allowbreak{}pass\_\allowbreak{}road\_\allowbreak{}LDV\_\allowbreak{}4W}
/
\texttt{Large\ Car\ and\ Truck}
& billion
passenger-km
& 35.54 &
18.13 \\
\texttt{trn\_\allowbreak{}freight\_\allowbreak{}road}
/
\texttt{Medium\ truck}
& billion
tonne-km &
6.587 &
3.587 \\
\texttt{trn\_\allowbreak{}freight}
/
\texttt{Freight\ Rail}
& billion
tonne-km &
4.404 &
0.644 \\
\texttt{trn\_\allowbreak{}freight}
/
\texttt{Domestic\ Ship}
& billion
tonne-km &
24.18 &
0 \\
\texttt{trn\_\allowbreak{}shipping\_\allowbreak{}intl}
/
\texttt{International\ Ship}
& billion
tonne-km &
2,837 &
0.058 \\
\end{longtable}

\textbf{Tags
with no
air-quality
species}
(greenhouse
gases,
fluorinated
gases or
hydrogen
only):

\begin{longtable}[]{@{}
  >{\raggedright\arraybackslash}p{(\columnwidth - 2\tabcolsep) * \real{0.5111}}
  >{\raggedright\arraybackslash}p{(\columnwidth - 2\tabcolsep) * \real{0.4889}}@{}}
\toprule\noalign{}
\begin{minipage}[b]{\linewidth}\raggedright
Sector /
subsector
\end{minipage}
&
\begin{minipage}[b]{\linewidth}\raggedright
Species
\end{minipage} \\
\midrule\noalign{}
\endhead
\bottomrule\noalign{}
\endlastfoot
\texttt{comm\ cooling},
\texttt{resid\ cooling\ modern\_\allowbreak{}d1}--\texttt{\_\allowbreak{}d10}
/
\texttt{electricity}
& HFC125,
HFC134a,
HFC143a,
HFC23,
HFC32 \\
\texttt{electricity\_\allowbreak{}net\_\allowbreak{}ownuse}
& SF6 \\
\texttt{industrial\ processes}
/
\texttt{Al\_\allowbreak{}Mg}
& SF6 \\
\texttt{industrial\ processes}
/
\texttt{HCFC\_\allowbreak{}22\_\allowbreak{}Prod}
& HFC23 \\
\texttt{industrial\ processes}
/
\texttt{adipic\ acid},
\texttt{nitric\ acid}
& N2O \\
\texttt{industrial\ processes}
/
\texttt{semiconductors}
& C2F6,
CF4,
HFC23,
SF6 \\
\texttt{urban\ processes}
/
\texttt{fire\_\allowbreak{}exting}
&
HFC227ea,
HFC236fa \\
\texttt{H2\ LDV},
\texttt{H2\ MHDV},
\texttt{LH2},
\texttt{H2\ liquid\ truck},
\texttt{H2\ pipeline}
(delivery)
& H2 \\
\texttt{biomass}
(bioenergy
crops) &
N2O\_AGR \\
\end{longtable}

Two
hydrofluorocarbon
(HFC)
sources,
\texttt{industrial\ processes\ /\allowbreak{}\ foams}
and
\texttt{urban\ processes\ /\allowbreak{}\ aerosols},
are active
but carry
no tag at
all.

\textbf{The
remaining
untagged
active
nodes} are
markets,
trade
nodes,
aggregators,
fuel
distribution,
water
supply,
labor and
capital,
renewable
and
nuclear
supply,
and
electricity-,
hydrogen-
or
human-powered
end uses.
The full
list, with
a reason
for each,
is in
\texttt{gcam\_\allowbreak{}kaist\_\allowbreak{}nonco2\_\allowbreak{}coverage\_\allowbreak{}by\_\allowbreak{}subsector.\allowbreak{}csv}.
Those
reasons
are our
screening
judgment
and have
not been
reviewed.

\hypertarget{b5.-where-base-year-values-come-from}{%
\subsubsection{B5.
Where
base-year
values
come
from}\label{b5.-where-base-year-values-come-from}}

\hypertarget{b5.1-two-ways-of-specifying-a-tag-in-stock-gcam}{%
\paragraph{B5.1
Two ways
of
specifying
a tag in
stock
GCAM}\label{b5.1-two-ways-of-specifying-a-tag-in-stock-gcam}}

In stock
GCAM v9's
inputs,
each
base-year
\texttt{Non-\allowbreak{}CO2}
entry is
specified
in one of
two ways:

\begin{itemize}
\tightlist
\item
  \texttt{\textless{}input-\allowbreak{}emissions\textgreater{}}:
  a mass.
  The
  coefficient
  is
  derived
  from it.
\item
  \texttt{\textless{}emiss-\allowbreak{}coef\textgreater{}}:
  a
  coefficient,
  given
  directly.
\end{itemize}

Each entry
also
declares a
driver:
\texttt{\textless{}input-\allowbreak{}driver\textgreater{}\textless{}input-\allowbreak{}name\textgreater{}…\textless{}/\allowbreak{}input-\allowbreak{}name\textgreater{}\textless{}/\allowbreak{}input-\allowbreak{}driver\textgreater{}}
for a
named
input, or
\texttt{\textless{}output-\allowbreak{}driver/\allowbreak{}\textgreater{}}
for the
technology's
output.
For South
Korea in
2015 and
2021:

\begin{longtable}[]{@{}
  >{\raggedright\arraybackslash}p{(\columnwidth - 4\tabcolsep) * \real{0.5349}}
  >{\raggedright\arraybackslash}p{(\columnwidth - 4\tabcolsep) * \real{0.2093}}
  >{\raggedright\arraybackslash}p{(\columnwidth - 4\tabcolsep) * \real{0.2558}}@{}}
\toprule\noalign{}
\begin{minipage}[b]{\linewidth}\raggedright
Sectors
\end{minipage}
&
\begin{minipage}[b]{\linewidth}\raggedright
Specified
as
\end{minipage}
&
\begin{minipage}[b]{\linewidth}\raggedright
Driver
\end{minipage} \\
\midrule\noalign{}
\endhead
\bottomrule\noalign{}
\endlastfoot
\texttt{electricity}
& mass &
output \\
\texttt{industrial\ processes},
\texttt{urban\ processes}
& mass &
output \\
Livestock
(\texttt{Beef},
\texttt{Dairy},
\texttt{Pork},
\texttt{Poultry},
\texttt{SheepGoat})
& mass &
output \\
Crops
(\texttt{\_\allowbreak{}AGR}
species) &
mass &
output \\
Crops
(\texttt{BC\_\allowbreak{}AWB},
\texttt{OC\_\allowbreak{}AWB})
&
coefficient
&
output \\
\texttt{UnmanagedLand}
& mass &
input
\texttt{land-\allowbreak{}input} \\
Buildings
(\texttt{comm\ *},
\texttt{resid\ *})
& mass &
named fuel
input \\
Transport
(\texttt{trn\_\allowbreak{}*})
& mass &
named fuel
input \\
Industrial
energy
(\texttt{chemical\ energy\ use},
\texttt{other\ industrial\ energy\ use},
\texttt{process\ heat\ *},
\texttt{ammonia},
\texttt{waste\ biomass\ for\ paper})
& mass &
named fuel
input \\
\texttt{agricultural\ energy\ use},
\texttt{construction\ energy\ use},
\texttt{mining\ energy\ use}
& mass &
named fuel
input \\
\texttt{iron\ and\ steel}
&
coefficient
&
output \\
Fuel
resources
(\texttt{coal},
\texttt{crude\ oil},
\texttt{natural\ gas})
&
coefficient
&
output \\
\end{longtable}

The other
\texttt{\_\allowbreak{}AWB}
species
that
appear in
GCAM-KAIST's
outputs
(for
example
\texttt{CO\_\allowbreak{}AWB},
\texttt{NOx\_\allowbreak{}AWB})
were not
found in
the stock
files
examined,
so how
they are
specified
is
unknown.

Stock
inputs
also
attach a
\texttt{gdp-\allowbreak{}control}
element
(with
\texttt{max-\allowbreak{}reduction}
and
\texttt{steepness})
to many
entries.
\texttt{emission\_\allowbreak{}factor\_\allowbreak{}controls.\allowbreak{}xml}
overrides
some
future
coefficients
directly.
For
example,
for South
Korea the
international
shipping
SO2
coefficient
is 0.16 in
2020,
falling to
0.13 by
2050, and
its GDP
control is
deleted.

\hypertarget{b5.2-gcam-kaist-compared-with-stock-gcam-v9}{%
\paragraph{B5.2
GCAM-KAIST
compared
with stock
GCAM
v9}\label{b5.2-gcam-kaist-compared-with-stock-gcam-v9}}

Extraction:
\texttt{results/\allowbreak{}diagnostics/\allowbreak{}capss\_\allowbreak{}calibration\_\allowbreak{}test/\allowbreak{}stock\_\allowbreak{}korea2.\allowbreak{}py}.

\textbf{Mass-specified
tags, 2015
and 2021,
AQ
species.}

\begin{longtable}[]{@{}
  >{\raggedright\arraybackslash}p{(\columnwidth - 2\tabcolsep) * \real{0.8000}}
  >{\raggedright\arraybackslash}p{(\columnwidth - 2\tabcolsep) * \real{0.2000}}@{}}
\toprule\noalign{}
\begin{minipage}[b]{\linewidth}\raggedright
\end{minipage}
&
\begin{minipage}[b]{\linewidth}\raggedright
Cells
\end{minipage} \\
\midrule\noalign{}
\endhead
\bottomrule\noalign{}
\endlastfoot
Present in
both,
identical
within
0.1\% &
2,244 \\
Present in
both,
different
& 7 \\
\textbf{Present
in both,
total} &
\textbf{2,251} \\
\end{longtable}

The seven
differing
cells are
\texttt{UnmanagedLand}
rows in
2015.
Beyond
them:

\begin{itemize}
\tightlist
\item
  \textbf{In
  stock,
  absent
  from
  GCAM-KAIST:}
  \texttt{chemical\ energy\ use\ /\allowbreak{}\ gas}
  in 2015,
  8.04 kt
  summed
  over AQ
  species,
  plus
  small
  agricultural-
  and
  mining-stationary
  entries.
\item
  \textbf{In
  GCAM-KAIST,
  absent
  from the
  stock
  files
  examined:}
  unmanaged-land fire subsectors (12.5 kt over 2015 and 2021).
\item
  \textbf{In GCAM-KAIST, carried forward from a stock 1975 coefficient:}
  \texttt{refining\ /\allowbreak{}\ biomass\ liquids\ /\allowbreak{}\ biodiesel}
  (1.40 kt over 2015 and 2021). Stock GCAM sets biodiesel only in 1975, as a
  coefficient per EJ of \texttt{regional\ biomassOil} input: 0.01 Tg/EJ for
  SO2, NOx and CO alike, and 0.0001 for N2O and NMVOC. GCAM-KAIST's 2015 and
  2021 biodiesel emissions equal that coefficient times the biodiesel's
  \texttt{regional\ biomassOil} input, to five significant figures.
\end{itemize}

\textbf{Coefficient-specified
tags.} For
\texttt{iron\ and\ steel}
and the
\texttt{coal},
\texttt{crude\ oil}
and
\texttt{natural\ gas}
resources,
GCAM-KAIST's
implied
coefficient
(emissions
÷ output)
equals the
stock
\texttt{emiss-\allowbreak{}coef}
to within
0.001\%.
This holds
for NOx,
SO2, CO
and NMVOC
in 2015
and 2021.

\textbf{Conclusion.}
In the
base
years,
GCAM-KAIST's
non-CO2
emissions
for South
Korea are,
with the
small
exceptions
listed,
stock GCAM
v9's.

\hypertarget{b5.3-upstream-of-stock-gcam-ceds-and-capss}{%
\paragraph{B5.3
Upstream
of stock
GCAM: CEDS
and
CAPSS}\label{b5.3-upstream-of-stock-gcam-ceds-and-capss}}

\begin{itemize}
\tightlist
\item
  GCAM's
  base-year
  non-CO2
  emissions
  are
  calibrated
  to CEDS.
  \texttt{ceds\_\allowbreak{}sector\_\allowbreak{}map.\allowbreak{}csv}
  maps
  CEDS
  sectors
  to GCAM
  sector
  groups.
  For
  example
  \texttt{1A1a\_\allowbreak{}Electricity-\allowbreak{}public}
  →
  \texttt{elec\_\allowbreak{}heat},
  and
  \texttt{1A1bc\_\allowbreak{}Other-\allowbreak{}transformation}
  →
  \texttt{industry\_\allowbreak{}energy}.
\item
  CEDS
  documentation
  states
  that its
  default
  estimates
  are
  scaled
  to
  country-level
  inventories
  where
  available.
  It lists
  \texttt{http:/\allowbreak{}/\allowbreak{}airemiss.\allowbreak{}nier.\allowbreak{}go.\allowbreak{}kr/\allowbreak{}},
  NIER's
  national
  air
  pollutant
  emission
  service,
  which
  publishes
  CAPSS,
  as the
  South
  Korea
  source
  (see
  Sources).
\item
  CEDS's own code confirms the Korea step (checked against the JGCRI/CEDS
  repository, March 2025 release). \texttt{code/\allowbreak{}module-E/\allowbreak{}E.South\_\allowbreak{}Korea\_\allowbreak{}emissions.R} reads the
  CAPSS inventory (\texttt{Korea\_\allowbreak{}CAPSS\_\allowbreak{}Emissions.xlsx}) for 1999--2021, and
  \texttt{code/\allowbreak{}module-F/\allowbreak{}F1.1.South\_\allowbreak{}Korea\_\allowbreak{}scaling.R} scales CEDS's default Korean
  estimates to it for SO2, NOx, CO, NMVOC and NH3 (1999--2021) and BC/OC (2011--2021). For
  BC/OC, \texttt{South\_\allowbreak{}Korea\_\allowbreak{}BCOC\_\allowbreak{}scaling\_\allowbreak{}method.csv} lists only \texttt{1A3b\_\allowbreak{}Road}.
\item
  \textbf{Not verified here:} which CEDS release GCAM v9 uses, and which Korean sectors and
  species that release actually scales to CAPSS. The current release also runs a second Korea
  scaling to EDGAR-HTAPv3.1 (\texttt{F1.1.South\_\allowbreak{}Korea\_\allowbreak{}EDGAR-HTAPv3.1\_\allowbreak{}scaling.R}; CO, NH3,
  NMVOC, NOx, SO2; 2000--2018), so in that release some Korean sector-years may end up matched to
  EDGAR-HTAP rather than directly to CAPSS.
\end{itemize}

\textbf{Why some cells match exactly and others are far off.} CEDS scales to CAPSS by broad
\emph{scaling sector}, not by detailed sector. \texttt{South\_\allowbreak{}Korea\_\allowbreak{}scaling\_\allowbreak{}mapping.csv} groups
CEDS's detailed sectors into groups that mirror CAPSS's categories: \texttt{energy} (public
electricity, heat, \texttt{1A1bc\_\allowbreak{}Other-transformation}), \texttt{industry} (industrial
fuel combustion), \texttt{1A3b\_\allowbreak{}Road}, \texttt{non road} (rail, shipping, domestic
aviation, construction machinery), \texttt{other\_\allowbreak{}combustion} (residential, commercial,
agriculture/forestry/fishing fuel), \texttt{fugitive}, \texttt{industry\_\allowbreak{}process} and
\texttt{solvent}. For each group, species and year, CEDS multiplies its own estimate so the
group total equals CAPSS. CEDS's split \emph{within} a group stays its own.

A CAPSS number therefore takes a three-step round trip before it is compared with CAPSS again:
CAPSS category $\to$ CEDS scaling sector $\to$ GCAM sector group (\texttt{ceds\_\allowbreak{}sector\_\allowbreak{}map.csv})
$\to$ CAPSS category (this repo's crosswalk). A cell comes back exact only when every step is
one-to-one for that category and the species was actually scaled:

\begin{itemize}
\tightlist
\item
  \textbf{Exact (B5.4):} agricultural NH3, road transport, industrial-process NH3 and waste SOx.
  Each passes through one CEDS group and one GCAM group that both map back to a single CAPSS
  category.
\item
  \textbf{Off because a step re-sorts mass:} CEDS counts \texttt{1A1bc\_\allowbreak{}Other-transformation}
  (which includes refining) under \texttt{energy}, but GCAM maps it to
  \texttt{industry\_\allowbreak{}energy}, so refinery emissions come back as Manufacturing industry rather
  than Energy production (B5.5; \textbf{Inferred}). Construction machinery is \texttt{non road}
  in CEDS but is split differently by GCAM and the crosswalk. Residential wood can land in
  Non-industry or Biomass burning.
\item
  \textbf{Off because the species was never scaled:} for BC, CEDS scales only road transport, so all
  other Korean BC keeps CEDS's default estimates. This is consistent with BC agreeing in Road
  transport (ratio 1.0) and differing 9--65$\times$ in the stationary-combustion categories (B5.4).
\item
  \textbf{Off because GCAM has no matching tag, or re-splits it:} mass on untagged technologies (B4),
  and placeholder-driven sources (B6.4) whose split differs from CAPSS's. This plausibly explains
  cells such as fuel storage and distribution VOCs at 12--18\% of CAPSS (\textbf{Inferred}).
\end{itemize}

The mapping search (\texttt{results/\allowbreak{}diagnostics/\allowbreak{}capss\_\allowbreak{}calibration\_\allowbreak{}test/\allowbreak{}}) shows the misses are
not a crosswalk-labelling problem: no assignment of the ambiguous GCAM sectors reproduces CAPSS. A
possible additional source of difference, not checked, is CAPSS revising past years after CEDS
ingested them.

\hypertarget{b5.4-agreement-with-capss}{%
\paragraph{B5.4
Agreement
with
CAPSS}\label{b5.4-agreement-with-capss}}

GCAM-KAIST
emissions
were
assigned
to CAPSS
first-level
categories
using the
repo
crosswalk
(\texttt{src/\allowbreak{}gcam\_\allowbreak{}emissions/\allowbreak{}reference/\allowbreak{}gcam\_\allowbreak{}kaist\_\allowbreak{}capss\_\allowbreak{}category\_\allowbreak{}crosswalk.\allowbreak{}csv})
plus the
two
overrides
in
\texttt{capss\_\allowbreak{}category\_\allowbreak{}comparison.\allowbreak{}py}:
\texttt{\_\allowbreak{}AWB}
→ Biomass
burning,
and
\texttt{industrial\ processes/\allowbreak{}solvents}
→ Solvent
use.
Eleven
category ×
pollutant
× year
cells
agree with
CAPSS to
within
0.12\%:

\begin{longtable}[]{@{}
  >{\raggedright\arraybackslash}p{(\columnwidth - 8\tabcolsep) * \real{0.2857}}
  >{\raggedright\arraybackslash}p{(\columnwidth - 8\tabcolsep) * \real{0.1786}}
  >{\raggedright\arraybackslash}p{(\columnwidth - 8\tabcolsep) * \real{0.1190}}
  >{\raggedright\arraybackslash}p{(\columnwidth - 8\tabcolsep) * \real{0.1786}}
  >{\raggedright\arraybackslash}p{(\columnwidth - 8\tabcolsep) * \real{0.2381}}@{}}
\toprule\noalign{}
\begin{minipage}[b]{\linewidth}\raggedright
CAPSS
category
\end{minipage}
&
\begin{minipage}[b]{\linewidth}\raggedright
Pollutant
\end{minipage}
&
\begin{minipage}[b]{\linewidth}\raggedright
Year
\end{minipage}
&
\begin{minipage}[b]{\linewidth}\raggedright
CAPSS (t)
\end{minipage}
&
\begin{minipage}[b]{\linewidth}\raggedright
GCAM-KAIST
(t)
\end{minipage} \\
\midrule\noalign{}
\endhead
\bottomrule\noalign{}
\endlastfoot
Agriculture
& NH3 &
2015 &
231,263 &
231,263.0 \\
Agriculture
& NH3 &
2021 &
200,384 &
200,384.0 \\
Road
transport
& NOx &
2021 &
287,279 &
287,340.6 \\
Road
transport
& SOx &
2015 & 209
& 208.8 \\
Road
transport
& SOx &
2021 & 248
& 248.0 \\
Road
transport
& NH3 &
2021 &
1,706 &
1,706.6 \\
Road
transport
& VOCs &
2015 &
46,145 &
46,146.7 \\
Industrial
process &
NH3 & 2015
& 39,432 &
39,467.3 \\
Industrial
process &
NH3 & 2021
& 41,953 &
41,991.1 \\
Waste
disposal &
SOx & 2015
& 2,119 &
2,119.1 \\
Waste
disposal &
SOx & 2021
& 1,382 &
1,383.5 \\
\end{longtable}

All eleven
cells come
from
mass-specified
tags
(B5.1).
The same
values are
present in
stock GCAM
v9 (B5.2).

Other
cells
agree much
less well.
BC, for
example,
differs by
large
factors in
the
stationary-combustion
categories
while
agreeing
closely in
Road
transport
and
Biomass
burning:

\begin{longtable}[]{@{}
  >{\raggedright\arraybackslash}p{(\columnwidth - 4\tabcolsep) * \real{0.3636}}
  >{\raggedright\arraybackslash}p{(\columnwidth - 4\tabcolsep) * \real{0.5065}}
  >{\raggedright\arraybackslash}p{(\columnwidth - 4\tabcolsep) * \real{0.1299}}@{}}
\toprule\noalign{}
\begin{minipage}[b]{\linewidth}\raggedright
CAPSS
category
\end{minipage}
&
\begin{minipage}[b]{\linewidth}\raggedright
BC ratio
GCAM-KAIST
/ CAPSS,
2015
\end{minipage}
&
\begin{minipage}[b]{\linewidth}\raggedright
2021
\end{minipage} \\
\midrule\noalign{}
\endhead
\bottomrule\noalign{}
\endlastfoot
Energy
production
& 8.9 &
11.1 \\
Manufacturing
industry &
10.1 &
36.8 \\
Non-industry
& 65.4 &
17.5 \\
Waste
disposal &
37.8 &
28.4 \\
Road
transport
& 1.0 &
1.0 \\
Biomass
burning &
1.0 &
0.9 \\
\end{longtable}

\textbf{Inferred.}
Agreement
to four or
five
significant
figures
cannot
arise from
independent
estimation.
Together
with B5.2
and B5.3,
it
indicates
that these
base-year
values are
CAPSS
totals,
passed
through
CEDS into
GCAM's
calibration
and
divided
among GCAM
technologies.

\hypertarget{b5.5-where-refinery-emissions-sit-inferred}{%
\paragraph{B5.5
Where
refinery
emissions
sit
(Inferred)}\label{b5.5-where-refinery-emissions-sit-inferred}}

\texttt{oil\ refining}
has no tag
(B4). CEDS
sector
\texttt{1A1bc\_\allowbreak{}Other-\allowbreak{}transformation},
which
includes
petroleum
refining,
is mapped
to GCAM's
\texttt{industry\_\allowbreak{}energy}
group
(B5.3).

The
comparison
tested
reassigning
\texttt{other\ industrial\ energy\ use\ /\allowbreak{}\ refined\ liquids}
from
Manufacturing
industry
to Energy
production
(CAPSS's
Energy
production
includes
refineries).
This was
done
together
with
moving
\texttt{refining\ /\allowbreak{}\ biomass\ liquids},
under 0.3
kt per
pollutant,
the other
way. The
result:

\begin{itemize}
\tightlist
\item
  Energy
  production
  CO and
  NH3 come
  within
  2.0\% of
  CAPSS in
  both
  years.
  NH3 in
  2015 is
  1,378.2
  t vs
  1,379 t.
\item
  Energy
  production
  VOCs
  come
  within
  2.3\%
  (2015)
  and
  6.1\%
  (2021).
\item
  Manufacturing
  industry
  CO in
  2021
  moves
  from
  85\%
  above
  CAPSS to
  3.6\%
  above.
\end{itemize}

This
indicates
that
refinery
fuel-combustion
emissions
are stored
in
\texttt{other\ industrial\ energy\ use\ /\allowbreak{}\ refined\ liquids},
not in
\texttt{refining}.

\hypertarget{b6.-implied-coefficients}{%
\subsubsection{B6.
Implied
coefficients}\label{b6.-implied-coefficients}}

\hypertarget{b6.1-method}{%
\paragraph{B6.1
Method}\label{b6.1-method}}

\texttt{results/\allowbreak{}diagnostics/\allowbreak{}nonco2\_\allowbreak{}coverage/\allowbreak{}driver\_\allowbreak{}consistent\_\allowbreak{}coef.\allowbreak{}py}
computes
each
coefficient
as
follows:

\begin{itemize}
\tightlist
\item
  \textbf{Numerator:}
  the
  tag's
  emissions,
  summed
  over
  vintages.
\item
  \textbf{Denominator:}
  the
  driver
  stock
  GCAM
  declares
  for that
  tag
  (B5.1).
  That is
  the
  named
  fuel
  input
  (\texttt{demand-\allowbreak{}physical})
  or the
  technology's
  output
  (\texttt{physical-\allowbreak{}output}).
\item
  \textbf{Included
  tags:}
  only
  mass-specified
  tags.
\item
  \textbf{Deciles:}
  income
  deciles
  \texttt{\_\allowbreak{}d2}--\texttt{\_\allowbreak{}d10}
  are
  dropped
  because
  they
  duplicate
  \texttt{\_\allowbreak{}d1}
  (shown
  below).
\item
  \textbf{Pollutants:}
  NOx,
  SO2,
  NMVOC,
  NH3, CO,
  BC and
  OC (no
  \texttt{\_\allowbreak{}AGR}
  or
  \texttt{\_\allowbreak{}AWB}).
\end{itemize}

\hypertarget{b6.2-distinct-technologies-with-the-same-coefficient}{%
\paragraph{B6.2
Distinct
technologies
with the
same
coefficient}\label{b6.2-distinct-technologies-with-the-same-coefficient}}

\textbf{Gas
power:
combined
cycle vs
steam/turbine.}

\begin{itemize}
\tightlist
\item
  Per unit
  of fuel
  burned,
  \texttt{gas\ (CC)}
  and
  \texttt{gas\ (steam/\allowbreak{}CT)}
  have
  identical
  coefficients
  in every
  base
  year
  from
  2005 to
  2021.
  The fuel
  burned
  is the
  gas
  input to
  the
  matching
  \texttt{elec\_\allowbreak{}*}
  sector.
\item
  The
  stock
  2021
  masses
  (0.006297
  and
  0.001779
  Tg NOx)
  are in
  the same
  ratio,
  3.54, as
  the two
  technologies'
  gas use
  (0.9943
  and
  0.2809
  EJ).
\item
  Per MWh
  generated,
  the
  coefficients
  differ
  only by
  the
  efficiency
  ratio:
  0.0404
  vs
  0.0605
  kg/MWh
  NOx in
  2021.
\item
  The two
  diverge
  only
  after
  2021.
\end{itemize}

\begin{longtable}[]{@{}
  >{\raggedright\arraybackslash}p{(\columnwidth - 12\tabcolsep) * \real{0.2747}}
  >{\raggedright\arraybackslash}p{(\columnwidth - 12\tabcolsep) * \real{0.1209}}
  >{\raggedright\arraybackslash}p{(\columnwidth - 12\tabcolsep) * \real{0.1209}}
  >{\raggedright\arraybackslash}p{(\columnwidth - 12\tabcolsep) * \real{0.1209}}
  >{\raggedright\arraybackslash}p{(\columnwidth - 12\tabcolsep) * \real{0.1209}}
  >{\raggedright\arraybackslash}p{(\columnwidth - 12\tabcolsep) * \real{0.1209}}
  >{\raggedright\arraybackslash}p{(\columnwidth - 12\tabcolsep) * \real{0.1209}}@{}}
\toprule\noalign{}
\begin{minipage}[b]{\linewidth}\raggedright
kg per TJ
of gas
\end{minipage}
&
\begin{minipage}[b]{\linewidth}\raggedright
2005
\end{minipage}
&
\begin{minipage}[b]{\linewidth}\raggedright
2010
\end{minipage}
&
\begin{minipage}[b]{\linewidth}\raggedright
2015
\end{minipage}
&
\begin{minipage}[b]{\linewidth}\raggedright
2021
\end{minipage}
&
\begin{minipage}[b]{\linewidth}\raggedright
2030
\end{minipage}
&
\begin{minipage}[b]{\linewidth}\raggedright
2050
\end{minipage} \\
\midrule\noalign{}
\endhead
\bottomrule\noalign{}
\endlastfoot
NOx,
\texttt{gas\ (CC)}
& 95.01 &
19.46 &
21.10 &
6.33 &
5.02 &
4.60 \\
NOx,
\texttt{gas\ (steam/\allowbreak{}CT)}
& 95.01 &
19.46 &
21.10 &
6.33 &
3.34 &
2.05 \\
CO,
\texttt{gas\ (CC)}
& 28.99 &
44.91 &
56.21 &
38.25 &
19.20 &
11.30 \\
CO,
\texttt{gas\ (steam/\allowbreak{}CT)}
& 28.99 &
44.91 &
56.21 &
38.25 &
20.17 &
12.04 \\
\end{longtable}

\textbf{Inferred.}
The
gas-power
total was
divided
between
the two
technologies
in
proportion
to fuel
use, not
according
to
technology-specific
emission
rates.

\textbf{Buildings.}

\begin{itemize}
\tightlist
\item
  \emph{Same
  coefficient
  across
  services.}
  In 2015
  and
  2021,
  \texttt{comm\ cooling},
  \texttt{comm\ heating}
  and
  \texttt{comm\ others}
  have the
  same
  coefficient
  per unit
  of gas
  for
  every
  pollutant
  (NOx in
  2021:
  67.45
  kg/TJ).
  \texttt{resid\ heating}
  and
  \texttt{resid\ others}
  match
  each
  other
  for
  coal,
  gas and
  refined
  liquids
  (checked
  in 2021;
  gas NOx:
  62.26
  kg/TJ).
\item
  \emph{Same
  coefficient
  across
  income
  deciles.}
  Across
  the ten
  income
  deciles,
  coefficients
  differ
  by at
  most
  0.09\%.
\end{itemize}

\textbf{Non-road
machinery:
equal
masses.}
In
\texttt{agricultural\ energy\ use},
\texttt{construction\ energy\ use}
and
\texttt{mining\ energy\ use},
the
\texttt{refined\ liquids}
technology
carries
\textbf{exactly
the same
mass} in
subsector
\texttt{mobile}
and
subsector
\texttt{stationary}.
This holds
in every
base year
(2005,
2010,
2015,
2021) and
for every
pollutant
checked
(NOx, CO,
SO2). The
equal
masses are
already
present in
stock GCAM
v9's
inputs
(checked
for
agricultural
and
construction
NOx in
2015 and
2021).

\begin{longtable}[]{@{}
  >{\raggedright\arraybackslash}p{(\columnwidth - 4\tabcolsep) * \real{0.4328}}
  >{\raggedright\arraybackslash}p{(\columnwidth - 4\tabcolsep) * \real{0.2537}}
  >{\raggedright\arraybackslash}p{(\columnwidth - 4\tabcolsep) * \real{0.3134}}@{}}
\toprule\noalign{}
\begin{minipage}[b]{\linewidth}\raggedright
2021, kt
\end{minipage}
&
\begin{minipage}[b]{\linewidth}\raggedright
NOx,
mobile
\end{minipage}
&
\begin{minipage}[b]{\linewidth}\raggedright
NOx,
stationary
\end{minipage} \\
\midrule\noalign{}
\endhead
\bottomrule\noalign{}
\endlastfoot
\texttt{agricultural\ energy\ use}
& 12.296 &
12.296 \\
\texttt{construction\ energy\ use}
& 31.103 &
31.103 \\
\texttt{mining\ energy\ use}
& 0.108 &
0.108 \\
\end{longtable}

\textbf{Steel:
EAF
routes.}
The stock
coefficients
themselves
give
electric-arc-furnace
steel made
from
direct-reduced
iron
(\texttt{EAF\ with\ DRI})
a far
higher NOx
coefficient
than EAF
steel made
from scrap
(\texttt{EAF\ with\ scrap}):

\begin{longtable}[]{@{}
  >{\raggedright\arraybackslash}p{(\columnwidth - 4\tabcolsep) * \real{0.5926}}
  >{\raggedright\arraybackslash}p{(\columnwidth - 4\tabcolsep) * \real{0.2037}}
  >{\raggedright\arraybackslash}p{(\columnwidth - 4\tabcolsep) * \real{0.2037}}@{}}
\toprule\noalign{}
\begin{minipage}[b]{\linewidth}\raggedright
NOx, kg
per tonne
of steel
\end{minipage}
&
\begin{minipage}[b]{\linewidth}\raggedright
2015
\end{minipage}
&
\begin{minipage}[b]{\linewidth}\raggedright
2021
\end{minipage} \\
\midrule\noalign{}
\endhead
\bottomrule\noalign{}
\endlastfoot
\texttt{EAF\ with\ DRI}
& 0.778 &
11.66 \\
\texttt{EAF\ with\ scrap}
& 0.134 &
0.216 \\
\texttt{BLASTFUR}
(blast
furnace) &
1.092 &
0.478 \\
\end{longtable}

In 2021
this gives
the two
EAF routes
nearly
equal NOx
masses
(4.67 and
4.75 kt),
although
their
outputs
are 0.400
and 21.98
Mt.

\hypertarget{b6.3-base-year-instability}{%
\paragraph{B6.3
Base-year
instability}\label{b6.3-base-year-instability}}

Across the
436
mass-specified
technology-pollutant
series
with a
positive
coefficient
in all
four of
2005,
2010, 2015
and 2021
(deciles
collapsed,
\texttt{ref}),
the ratio
of the
largest to
the
smallest
coefficient
is:

\begin{longtable}[]{@{}
  >{\raggedright\arraybackslash}p{(\columnwidth - 2\tabcolsep) * \real{0.3235}}
  >{\raggedright\arraybackslash}p{(\columnwidth - 2\tabcolsep) * \real{0.6765}}@{}}
\toprule\noalign{}
\begin{minipage}[b]{\linewidth}\raggedright
Statistic
\end{minipage}
&
\begin{minipage}[b]{\linewidth}\raggedright
Value
\end{minipage} \\
\midrule\noalign{}
\endhead
\bottomrule\noalign{}
\endlastfoot
Median &
2.04 \\
Share
above 1.5×
& 69\% \\
Share
above 2× &
52\% \\
90th
percentile
& 11.6 \\
Maximum &
134.7 (OC
from
residential
coal
heating) \\
\end{longtable}

Examples:

\begin{longtable}[]{@{}
  >{\raggedright\arraybackslash}p{(\columnwidth - 8\tabcolsep) * \real{0.4894}}
  >{\raggedright\arraybackslash}p{(\columnwidth - 8\tabcolsep) * \real{0.1277}}
  >{\raggedright\arraybackslash}p{(\columnwidth - 8\tabcolsep) * \real{0.1277}}
  >{\raggedright\arraybackslash}p{(\columnwidth - 8\tabcolsep) * \real{0.1277}}
  >{\raggedright\arraybackslash}p{(\columnwidth - 8\tabcolsep) * \real{0.1277}}@{}}
\toprule\noalign{}
\begin{minipage}[b]{\linewidth}\raggedright
Coefficient
\end{minipage}
&
\begin{minipage}[b]{\linewidth}\raggedright
2005
\end{minipage}
&
\begin{minipage}[b]{\linewidth}\raggedright
2010
\end{minipage}
&
\begin{minipage}[b]{\linewidth}\raggedright
2015
\end{minipage}
&
\begin{minipage}[b]{\linewidth}\raggedright
2021
\end{minipage} \\
\midrule\noalign{}
\endhead
\bottomrule\noalign{}
\endlastfoot
Coal power
(\texttt{coal\ (conv\ pul)})
NOx,
kg/MWh &
1.753 &
0.509 &
0.458 &
0.150 \\
Coal power
SO2,
kg/MWh &
0.408 &
0.214 &
0.267 &
0.071 \\
Gas
combined-cycle
(\texttt{gas\ (CC)})
NOx,
kg/MWh &
0.684 &
0.124 &
0.134 &
0.040 \\
Oil power
(\texttt{refined\ liquids\ (steam/\allowbreak{}CT)})
SO2,
kg/MWh &
0.395 &
0.424 &
0.521 &
0.374 \\
Cement-kiln
coal NOx,
kg/GJ coal
& 0.318 &
0.346 &
0.247 &
0.264 \\
Commercial
gas
(\texttt{comm\ others})
NOx, kg/GJ
gas &
0.0366 &
0.0402 &
0.0416 &
0.0674 \\
Car
petrol/diesel
(\texttt{Liquids})
NOx, kg/GJ
fuel &
0.229 &
0.117 &
0.119 &
0.072 \\
\end{longtable}

\textbf{Inferred.}
For
mass-specified
tags, each
base-year
coefficient
equals the
inventory
mass
divided by
GCAM's
modelled
activity,
so any
difference
between
GCAM's
activity
and the
activity
behind the
inventory
is
absorbed
into the
coefficient.

\hypertarget{b6.4-placeholder-driven-tags}{%
\paragraph{B6.4
Placeholder-driven
tags}\label{b6.4-placeholder-driven-tags}}

Every
tagged
subsector
of
\texttt{industrial\ processes}
and
\texttt{urban\ processes}
takes
\texttt{misc\ emissions\ sources},
an
\texttt{unlimited-\allowbreak{}resource},
as input.
Stock GCAM
declares
an output
driver for
these
tags. That
output is
a fixed
placeholder
(unit
\texttt{NA}),
not a
production
quantity:

\begin{itemize}
\tightlist
\item
  In 2021
  all
  eight
  \texttt{industrial\ processes}
  subsectors
  output
  0.000243714.
\item
  The
  \texttt{urban\ processes}
  subsectors
  output
  0.0008.
\end{itemize}

The NOx
and SO2 on
\texttt{industrial\ processes\ /\allowbreak{}\ other\ industrial\ processes}
fall from
2021 to
2050
almost
identically
in both
scenarios:

\begin{longtable}[]{@{}
  >{\raggedright\arraybackslash}p{(\columnwidth - 6\tabcolsep) * \real{0.4571}}
  >{\raggedright\arraybackslash}p{(\columnwidth - 6\tabcolsep) * \real{0.1571}}
  >{\raggedright\arraybackslash}p{(\columnwidth - 6\tabcolsep) * \real{0.2000}}
  >{\raggedright\arraybackslash}p{(\columnwidth - 6\tabcolsep) * \real{0.1857}}@{}}
\toprule\noalign{}
\begin{minipage}[b]{\linewidth}\raggedright
\texttt{other\ industrial\ processes}
\end{minipage}
&
\begin{minipage}[b]{\linewidth}\raggedright
2021
\end{minipage}
&
\begin{minipage}[b]{\linewidth}\raggedright
2050
\texttt{ref}
\end{minipage}
&
\begin{minipage}[b]{\linewidth}\raggedright
2050
\texttt{nz}
\end{minipage} \\
\midrule\noalign{}
\endhead
\bottomrule\noalign{}
\endlastfoot
NOx, kt &
39.10 &
13.62 &
13.36 \\
SO2, kt &
83.81 &
29.19 &
28.64 \\
\end{longtable}

The table
below
gives the
share of
GCAM-KAIST's
2021
Korean
emissions
on
placeholder-driven
tags. It
uses
\texttt{ref}
and
excludes
\texttt{trn\_\allowbreak{}aviation\_\allowbreak{}intl}
and
\texttt{trn\_\allowbreak{}shipping\_\allowbreak{}intl}.

\begin{longtable}[]{@{}
  >{\raggedright\arraybackslash}p{(\columnwidth - 4\tabcolsep) * \real{0.2239}}
  >{\raggedright\arraybackslash}p{(\columnwidth - 4\tabcolsep) * \real{0.4776}}
  >{\raggedright\arraybackslash}p{(\columnwidth - 4\tabcolsep) * \real{0.2985}}@{}}
\toprule\noalign{}
\begin{minipage}[b]{\linewidth}\raggedright
Pollutant
\end{minipage}
&
\begin{minipage}[b]{\linewidth}\raggedright
On
placeholder-driven
tags
\end{minipage}
&
\begin{minipage}[b]{\linewidth}\raggedright
Share of
total
\end{minipage} \\
\midrule\noalign{}
\endhead
\bottomrule\noalign{}
\endlastfoot
VOCs & 760
kt &
81\% \\
SO2 & 85
kt &
55\% \\
NH3 & 42
kt &
17\% \\
NOx & 50
kt &
6\% \\
CO & 35 kt
& 6\% \\
\end{longtable}

The VOCs
are mainly
\texttt{solvents}
(574 kt),
\texttt{other\ industrial\ processes}
(128 kt)
and
\texttt{landfills}
(48 kt).

\hypertarget{b6.5-future-years}{%
\paragraph{B6.5
Future
years}\label{b6.5-future-years}}

Of 463
mass-specified
series
with a
positive
2021
coefficient
(deciles
collapsed,
\texttt{ref}),
comparing
the 2050
coefficient
with the
2021 one:

\begin{longtable}[]{@{}
  >{\raggedright\arraybackslash}p{(\columnwidth - 4\tabcolsep) * \real{0.5208}}
  >{\raggedright\arraybackslash}p{(\columnwidth - 4\tabcolsep) * \real{0.2500}}
  >{\raggedright\arraybackslash}p{(\columnwidth - 4\tabcolsep) * \real{0.2292}}@{}}
\toprule\noalign{}
\begin{minipage}[b]{\linewidth}\raggedright
Behaviour
\end{minipage}
&
\begin{minipage}[b]{\linewidth}\raggedright
Series
\end{minipage}
&
\begin{minipage}[b]{\linewidth}\raggedright
Share
\end{minipage} \\
\midrule\noalign{}
\endhead
\bottomrule\noalign{}
\endlastfoot
Frozen
(within
±1\%) &
213 &
46\% \\
Declining
& 243 &
52\% \\
Rising & 7
& 2\% \\
\end{longtable}

Frozen
series are
most
common in
\texttt{process\ heat\ cement},
\texttt{construction\ energy\ use},
\texttt{UnmanagedLand}
and
residential
heating.
Declining
series are
most
common in
\texttt{electricity},
\texttt{chemical\ energy\ use},
\texttt{other\ industrial\ energy\ use}
and cars.

In 2050,
409 of 439
coefficients
(93\%) are
identical
in
\texttt{nz}
and
\texttt{ref}.

\begin{longtable}[]{@{}
  >{\raggedright\arraybackslash}p{(\columnwidth - 8\tabcolsep) * \real{0.4000}}
  >{\raggedright\arraybackslash}p{(\columnwidth - 8\tabcolsep) * \real{0.1412}}
  >{\raggedright\arraybackslash}p{(\columnwidth - 8\tabcolsep) * \real{0.1412}}
  >{\raggedright\arraybackslash}p{(\columnwidth - 8\tabcolsep) * \real{0.1412}}
  >{\raggedright\arraybackslash}p{(\columnwidth - 8\tabcolsep) * \real{0.1765}}@{}}
\toprule\noalign{}
\begin{minipage}[b]{\linewidth}\raggedright
Coefficient
\end{minipage}
&
\begin{minipage}[b]{\linewidth}\raggedright
2021
\end{minipage}
&
\begin{minipage}[b]{\linewidth}\raggedright
2030
\end{minipage}
&
\begin{minipage}[b]{\linewidth}\raggedright
2050
\end{minipage}
&
\begin{minipage}[b]{\linewidth}\raggedright
Behaviour
\end{minipage} \\
\midrule\noalign{}
\endhead
\bottomrule\noalign{}
\endlastfoot
Cement-kiln
coal NOx,
kg/GJ &
0.264 &
0.264 &
0.264 &
frozen \\
Commercial
gas NOx,
kg/GJ &
0.0674 &
0.0674 &
0.0674 &
frozen \\
Oil power
NOx,
kg/MWh &
0.320 &
0.320 &
0.320 &
frozen \\
Oil power
SO2,
kg/MWh &
0.374 &
0.193 &
0.108 &
declining \\
Coal power
NOx,
kg/MWh &
0.150 &
0.077 &
0.043 &
declining \\
Coal power
SO2,
kg/MWh &
0.071 &
0.037 &
0.021 &
declining \\
Car
petrol/diesel
NOx, kg/GJ
& 0.072 &
0.037 &
0.021 &
declining \\
\end{longtable}

The stock
inputs
contain
the
mechanisms
that
produce
this
behaviour
(\texttt{gdp-\allowbreak{}control},
and future
\texttt{emiss-\allowbreak{}coef}
overrides
in
\texttt{emission\_\allowbreak{}factor\_\allowbreak{}controls.\allowbreak{}xml}).
For South Korea, that file sets 2025 coefficients for the six CCS power
technologies (for example \texttt{gas\ (CC\ CCS)} NOx: 0.00246) and SO2
coefficients for domestic and international shipping.
GCAM-KAIST's own control settings
cannot be
inspected,
because
its input
files are
not
available.

\hypertarget{b7.-summary-what-the-encoding-can-and-cannot-represent}{%
\subsubsection{B7.
Summary:
what the
encoding
can and
cannot
represent}\label{b7.-summary-what-the-encoding-can-and-cannot-represent}}

\begin{longtable}[]{@{}
  >{\raggedright\arraybackslash}p{(\columnwidth - 2\tabcolsep) * \real{0.5000}}
  >{\raggedright\arraybackslash}p{(\columnwidth - 2\tabcolsep) * \real{0.5000}}@{}}
\toprule\noalign{}
\begin{minipage}[b]{\linewidth}\raggedright
Change in
a scenario
\end{minipage}
&
\begin{minipage}[b]{\linewidth}\raggedright
How
GCAM-KAIST's
non-CO2
emissions
respond
\end{minipage} \\
\midrule\noalign{}
\endhead
\bottomrule\noalign{}
\endlastfoot
Shift
between
tagged
technologies
(for
example
coal → gas
power) &
Emissions
change at
each
technology's
coefficient. \\
Shift
between
technologies
that share
a
coefficient
(gas CC ↔
steam/CT
per unit
of fuel;
heating ↔
cooling ↔
other;
deciles) &
Base-year
coefficients
do not
distinguish
them. \\
Shift to
an
untagged
technology
(hybrid
vehicles,
hydrogen-based
DRI, steel
or cement
with CCS)
&
Emissions
of that
activity
are
zero. \\
Change in
industrial
output or
refinery
throughput
& Process
emissions
do not
respond:
they are
absent, or
sit on a
placeholder
driver.
Refinery
combustion
follows
\texttt{other\ industrial\ energy\ use\ /\allowbreak{}\ refined\ liquids}
(Inferred,
B5.5). \\
Pollution-control
policy
(scrubbers,
emission
standards)
& No
lever.
Coefficients
are
frozen, or
follow a
curve
shared by
both
scenarios. \\
Primary
PM2.5 /
PM10 & Not
represented. \\
\end{longtable}

\hypertarget{b8.-limitations-of-this-audit}{%
\subsubsection{B8.
Limitations
of this
audit}\label{b8.-limitations-of-this-audit}}

\begin{itemize}
\tightlist
\item
  \textbf{KAIST
  inputs.}
  GCAM-KAIST's
  input
  XMLs
  were not
  examined.
  Statements
  about
  how its
  tags are
  set rely
  on its
  outputs
  and on
  stock
  GCAM
  v9's
  inputs.
\item
  \textbf{CEDS
  link.}
  The CEDS
  → CAPSS
  link
  rests on
  CEDS
  documentation
  and on
  the
  numerical
  agreement
  in B5.4.
  It has
  not been
  traced
  through
  the CEDS
  release
  used by
  GCAM v9.
\item
  \textbf{Refinery
  location.}
  B5.5 is
  an
  inference
  from a
  sector
  map and
  a
  mapping
  test,
  not from
  an
  explicit
  source
  record.
\item
  \textbf{Coefficient
  scope.}
  Implied
  coefficients
  use the
  driver
  stock
  GCAM
  declares.
  Where
  GCAM-KAIST
  declares
  a
  different
  driver,
  its true
  coefficient
  would
  differ.
\item
  \textbf{Reasons
  for
  untagged
  nodes.}
  The
  classification
  of
  untagged
  active
  nodes
  (B4) is
  our
  screening
  judgment
  and has
  not been
  reviewed.
\item
  \textbf{Region.}
  Only
  South
  Korea
  was
  examined.
\end{itemize}

\hypertarget{b9.-reproducing-the-evidence}{%
\subsubsection{B9.
Reproducing
the
evidence}\label{b9.-reproducing-the-evidence}}

All
scripts
and
outputs
below are
local
diagnostics
under
\texttt{results/\allowbreak{}},
which is
untracked.
They
should be
moved into
the
package
before
this
document
is
circulated
as a
reproducible
record.

\begin{longtable}[]{@{}
  >{\raggedright\arraybackslash}p{(\columnwidth - 4\tabcolsep) * \real{0.3333}}
  >{\raggedright\arraybackslash}p{(\columnwidth - 4\tabcolsep) * \real{0.3333}}
  >{\raggedright\arraybackslash}p{(\columnwidth - 4\tabcolsep) * \real{0.3333}}@{}}
\toprule\noalign{}
\begin{minipage}[b]{\linewidth}\raggedright
Evidence
\end{minipage}
&
\begin{minipage}[b]{\linewidth}\raggedright
Script
\end{minipage}
&
\begin{minipage}[b]{\linewidth}\raggedright
Output
\end{minipage} \\
\midrule\noalign{}
\endhead
\bottomrule\noalign{}
\endlastfoot
Native
emissions
detail
(all tags)
&
\texttt{src/\allowbreak{}gcam\_\allowbreak{}emissions/\allowbreak{}gcam/\allowbreak{}nonco2\_\allowbreak{}by\_\allowbreak{}sector.\allowbreak{}py}
+
\texttt{queries/\allowbreak{}nonco2\_\allowbreak{}emissions\_\allowbreak{}by\_\allowbreak{}sector.\allowbreak{}xq}
&
\texttt{data/\allowbreak{}interim/\allowbreak{}gcam\_\allowbreak{}native\_\allowbreak{}emissions\_\allowbreak{}by\_\allowbreak{}sector/\allowbreak{}gcam\_\allowbreak{}kaist\_\allowbreak{}nonco2\_\allowbreak{}emissions\_\allowbreak{}detail.\allowbreak{}csv} \\
Activity
and tag
coverage
(B4) &
\texttt{results/\allowbreak{}diagnostics/\allowbreak{}nonco2\_\allowbreak{}coverage/\allowbreak{}activity\_\allowbreak{}inventory.\allowbreak{}xq}
(run with
\texttt{basex\ -\allowbreak{}b\ db=\allowbreak{}ref},
then
\texttt{db=\allowbreak{}nz}),
\texttt{build\_\allowbreak{}coverage.\allowbreak{}py}
&
\texttt{gcam\_\allowbreak{}kaist\_\allowbreak{}nonco2\_\allowbreak{}coverage\_\allowbreak{}by\_\allowbreak{}technology.\allowbreak{}csv},
\texttt{.\allowbreak{}.\allowbreak{}.\allowbreak{}\_\allowbreak{}by\_\allowbreak{}subsector.\allowbreak{}csv} \\
Activity
drivers by
technology
&
\texttt{results/\allowbreak{}diagnostics/\allowbreak{}policy\_\allowbreak{}experiment/\allowbreak{}drivers.\allowbreak{}xq}
&
\texttt{drv\_\allowbreak{}ref.\allowbreak{}csv},
\texttt{drv\_\allowbreak{}nz.\allowbreak{}csv} \\
Power fuel
use &
\texttt{results/\allowbreak{}diagnostics/\allowbreak{}nonco2\_\allowbreak{}coverage/\allowbreak{}elecfuel.\allowbreak{}xq}
&
\texttt{elecfuel\_\allowbreak{}ref.\allowbreak{}csv},
\texttt{elecfuel\_\allowbreak{}nz.\allowbreak{}csv} \\
Stock GCAM
v9 Korea
specifications
(B5.1--B5.2)
&
\texttt{results/\allowbreak{}diagnostics/\allowbreak{}capss\_\allowbreak{}calibration\_\allowbreak{}test/\allowbreak{}stock\_\allowbreak{}korea2.\allowbreak{}py}
&
\texttt{stock\_\allowbreak{}korea2.\allowbreak{}csv} \\
CAPSS
agreement
(B5.4--B5.5)
&
\texttt{results/\allowbreak{}diagnostics/\allowbreak{}capss\_\allowbreak{}calibration\_\allowbreak{}test/\allowbreak{}mapping\_\allowbreak{}search.\allowbreak{}py}
&
\texttt{capss\_\allowbreak{}mapping\_\allowbreak{}search\_\allowbreak{}default.\allowbreak{}csv},
\texttt{.\allowbreak{}.\allowbreak{}.\allowbreak{}\_\allowbreak{}best.\allowbreak{}csv} \\
Driver-consistent
coefficients
(B6.1,
B6.3,
B6.5) &
\texttt{results/\allowbreak{}diagnostics/\allowbreak{}nonco2\_\allowbreak{}coverage/\allowbreak{}driver\_\allowbreak{}consistent\_\allowbreak{}coef.\allowbreak{}py}
&
\texttt{driver\_\allowbreak{}consistent\_\allowbreak{}coef.\allowbreak{}csv} \\
Per-fuel
power
coefficients
and
placeholder
shares
(B6.2,
B6.4) &
\texttt{results/\allowbreak{}diagnostics/\allowbreak{}nonco2\_\allowbreak{}coverage/\allowbreak{}implied\_\allowbreak{}ef.\allowbreak{}py}
(its
sections
D--E are
superseded
by the
line
above) &
\texttt{implied\_\allowbreak{}ef\_\allowbreak{}long.\allowbreak{}csv} \\
\end{longtable}

Related
analyses
that build
on this
audit, not
part of
it:

\begin{itemize}
\tightlist
\item
  the
  CAPSS
  mapping
  review
  (\texttt{results/\allowbreak{}diagnostics/\allowbreak{}nonco2\_\allowbreak{}coverage/\allowbreak{}gcam\_\allowbreak{}kaist\_\allowbreak{}capss\_\allowbreak{}mapping\_\allowbreak{}review.\allowbreak{}csv});
\item
  the
  mapping
  search
  (\texttt{results/\allowbreak{}diagnostics/\allowbreak{}capss\_\allowbreak{}calibration\_\allowbreak{}test/\allowbreak{}});
\item
  the
  policy-sensitivity
  experiment
  (\texttt{results/\allowbreak{}diagnostics/\allowbreak{}policy\_\allowbreak{}experiment/\allowbreak{}}).
\end{itemize}

\hypertarget{sources}{%
\subsubsection{Sources}\label{sources}}

\begin{itemize}
\tightlist
\item
  JGCRI,
  \emph{Community
  Emissions
  Data
  System
  (CEDS)},
  \url{https://www.github.com/JGCRI/CEDS}.
\item
  CEDS
  reviewer-response
  supplement,
  ESSD
  preprint
  essd-2020-103,
  \url{https://essd.copernicus.org/preprints/essd-2020-103/essd-2020-103-AC1-supplement.pdf}.
\item
  Stock
  GCAM v9
  files
  listed
  in B1.
\end{itemize}
